\subsection{Morfismos de tipo finito.}
\label{subsection:IV.1.5-fr}

\begin{proposition}[1.5.1]
\label{IV.1.5.1-fr}
Sean $f:X\to Y$ un morfismo, $(U_\alpha)$ un recubrimiento de $Y$ formado
por abiertos afines. Las condiciones siguientes son equivalentes:
\begin{enumerate}
  \item[\textup{a)}] $f$ es localmente de tipo finito y cuasicompacto.
  \item[\textup{b)}] Para todo $\alpha$, $f^{-1}(U_\alpha)$ es unión finita de abiertos
  afines $V_{\alpha i}$ tales que cada anillo $A(V_{\alpha i})$ es una
  $A(U_\alpha)$-álgebra de tipo finito.
  \item[\textup{c)}] Para todo abierto afín $U$ de $Y$, $f^{-1}(U)$ es unión finita
  de abiertos afines $V_j$ tales que los anillos $A(V_j)$ son
  $A(U)$-álgebras de tipo finito.
\end{enumerate}
\end{proposition}

Para la demostración, véase (\textbf{I}, 6.3.2 y 6.6.3).

\begin{definition}[1.5.2]
\label{IV.1.5.2-fr}
Se dice que un morfismo $f$ es de tipo finito si verifica las condiciones
equivalentes de la proposición \textup{(1.5.1)}.
\end{definition}

\oldpage[IV-1]{234}

\begin{proposition}[1.5.3]
\label{IV.1.5.3-fr}
Sea $f:X\to Y$ un morfismo de tipo finito; si $Y$ es noetheriano, lo mismo
ocurre con $X$.
\end{proposition}

Para la demostración, véase (\textbf{I}, 6.3.7).

\begin{proposition}[1.5.4]
\label{IV.1.5.4-fr}
\begin{enumerate}
  \item[\textup{(i)}] Sea $j:X\to Y$ un morfismo de inmersión. Si $j$ es
  cuasicompacto (lo que ocurre si $j$ es una inmersión cerrada, o si
  el espacio subyacente a $X$ es noetheriano, o si el espacio subyacente a $Y$
  es localmente noetheriano), $j$ es de tipo finito.
  \item[\textup{(ii)}] El compuesto de dos morfismos de tipo finito es de tipo finito.
  \item[\textup{(iii)}] Si $f:X\to Y$ es un $S$-morfismo de tipo finito,
  $f_{(S')}:X_{(S')}\to Y_{(S')}$ es de tipo finito para toda extensión
  $S'\to S$ del preesquema de base.
  \item[\textup{(iv)}] Si $f:X\to X'$ y $g:Y\to Y'$ son dos $S$-morfismos de tipo
  finito, $f\times_Sg$ es de tipo finito.
  \item[\textup{(v)}] Si el compuesto $g\circ f$ de dos morfismos es de tipo finito y
  si $g$ es cuasiseparado o si $X$ es noetheriano, $f$ es de tipo finito.
  \item[\textup{(vi)}] Si un morfismo $f$ es de tipo finito, lo mismo ocurre con
  $f_{\mathrm{red}}$.
\end{enumerate}
\end{proposition}

Esto resulta inmediatamente (salvo el caso (v) en que $X$ es noetheriano, demostrado en
(\textbf{I}, 6.3.6)) de las definiciones y de (1.1.2), (1.2.4) y (1.3.4).

\begin{corollary}[1.5.5]
\label{IV.1.5.5-fr}
Sea $f:X\to Y$ un morfismo de tipo finito. Para todo morfismo $Y'\to Y$ tal
que $Y'$ sea noetheriano, $X\times_YY'$ es noetheriano.
\end{corollary}

Esto resulta de (1.5.4, (iii)) y de (1.5.3).

\begin{corollary}[1.5.6]
\label{IV.1.5.6-fr}
Sea $X$ un preesquema de tipo finito sobre un preesquema localmente noetheriano $S$.
Entonces todo $S$-morfismo $f:X\to Y$ es de tipo finito.
\end{corollary}

Para la demostración, véase (\textbf{I}, 6.3.9).

\begin{proposition}[1.5.7]
\label{IV.1.5.7-fr}
Sea $\varphi:A\to B$ un homomorfismo de anillos. Para que el morfismo
correspondiente $f:\operatorname{Spec}(B)\to\operatorname{Spec}(A)$ sea de tipo finito, es necesario y suficiente
que $\varphi$ haga de $B$ una $A$-álgebra de tipo finito.
\end{proposition}

Para la demostración, véase (\textbf{I}, 6.3.3).

